\documentclass[10pt,a4paper]{amsart}
\usepackage[margin=19mm]{geometry}
\usepackage{amsmath,amssymb,mathtools}
\usepackage[scr=rsfs,cal=euler]{mathalfa}
\usepackage{xfrac}
\usepackage[unicode,hidelinks,bookmarks=false]{hyperref}
\newcommand{\ie}{i.e.\ }
\newcommand{\cf}{cf.\ }
\newcommand{\resp}{respectively\ }
\newcommand{\Subsection}{\subsection}
\providecommand{\nobreakdash}{\nobreak\hbox{-}}
\setlength{\emergencystretch}{2em}
\multlinegap0pt
\usepackage{bm}
\usepackage{bbm}
\usepackage{dsfont}
\usepackage{xspace}
\usepackage{cancel}
\usepackage{mathtools}
\usepackage{amsthm}
\makeatletter
\@ifundefined{theorem}{\newtheorem{theorem}{Theorem}}{}
\@ifundefined{lemma}{\newtheorem{lemma}[theorem]{Lemma}}{}
\makeatother
\newcommand{\forceP}{\mathbb{P}}
\newcommand{\forceQ}{\mathbb{Q}}
\title[A Universe with a $\Delta^1\_n$-definable well-order of the r]{A Universe with a $\Delta^1\_n$-definable well-order of the reals, $\mathsf{CH}$ and $\Pi^1\_n$-Uniformization}
\author{Stefan Hoffelner}
\date{Source: arXiv:2506.21778v1 (CC BY 4.0)}
\begin{document}
\maketitle

\noindent\emph{Source and licence.} A Universe with a $\Delta^1_n$-definable well-order of the reals, $\mathsf{CH}$ and $\Pi^1_n$-Uniformization. Version arXiv:2506.21778v1 (CC BY 4.0), distributed under the Creative Commons Attribution 4.0 International licence, as recorded in the arXiv licence metadata for this exact version and shown on the abstract page https://arxiv.org/abs/2506.21778v1. Full rights notice: licenses/arxiv-2506.21778-CC-BY-4.0.txt.\par\smallskip

\noindent\textit{Editorial scope.} Counted unit: the Lemma whose source label is \texttt{nounwantedcodes} in the article (source0.tex lines 620--624, source label \texttt{nounwantedcodes}), delivered together with the article's own complete original proof of it (source0.tex lines 625--682, one \texttt{proof} environment of 58 lines). No mathematics is written, repaired or re-derived: statement and proof are the article's own text line for line. Required context retained uncounted: PullingOutATree (lines 401--408). Every definition, display and label of those lines is retained unchanged. Omitted from this package and not redistributed: 1--400, 409--619, 683--1126. Nothing inside a retained excerpt is omitted or altered: every retained body is delivered exactly as the article's lines are written, so no omission receipt of any kind accompanies this package. Citations: the retained excerpts carry no citation command at all, so no bibliography entry of the article is transcribed and no reference is redistributed. Reference adaptation: none was needed and none was made; every reference inside the delivered unit resolves inside it, so no unresolved reference or citation is produced. Printed slips kept literal: no printed word, symbol, hypothesis or number of the counted statement or of its proof was altered. Layout and class: the article's own class is \texttt{article} with packages including inputenc, boxedminipage, amsfonts, amsmath, amssymb, graphicx, amsthm, t1enc, subfig, cancel, textcmds, mathabx; the wrapper is the lane's pinned \texttt{amsart} editorial wrapper at 10pt on A4, which restates the theorem-like environments the retained text needs and adds the article's own macro definitions that the retained text needs (\texttt{\textbackslash forceP}, \texttt{\textbackslash forceQ}), with extra wrapper packages (bm, bbm, dsfont, xspace, cancel, mathtools). The delivered numbering is the unit's own. Assets: no retained span contains a figure environment, a graphics-inclusion command, a PDF-page inclusion, a table or a TikZ picture; no image file is copied into this package, the package declares no asset dependency and no image file is redistributed with it. Licence: version arXiv:2506.21778v1 of this article is distributed under the Creative Commons Attribution 4.0 International licence, as recorded in the arXiv licence metadata for this exact version and shown on the abstract page https://arxiv.org/abs/2506.21778v1. A full rights notice is delivered at licenses/arxiv-2506.21778-CC-BY-4.0.txt. Characters: every retained excerpt is pure ASCII, so the delivered engine, XeLaTeX with its default Latin Modern text font, reports no missing character.
\begin{lemma}\label{PullingOutATree}
Let $\vec{S} \ast \vec{b} \times G$ be an $L$-generic filter for $\forceQ^0 \ast \forceQ^1 \times \forceQ^2$.
Let $A \subset \omega_1$, $A \in L$, and let $\vec{b} \upharpoonright A = (b_{\alpha} \mid \alpha \in A)$.
Work in $W':=L[\vec{S} \ast (\vec{b} \upharpoonright A) \times G]$ and let $\forceP = (\forceP_{\alpha}, \forceP(\alpha) \mid \alpha < \delta < \omega_1) \in W' $ be 
a countable support iteration such that for every factor $\forceP(\alpha)$, either $\forceP(\alpha)$ is a $\forceP_{\alpha}$-name of $S_{-1}$, or $\forceP(\alpha)$ is a $\forceP_{\alpha}$-name of a proper, $\aleph_1$-sized forcing which preserves Suslin trees.

Let $\beta \notin A$. Then $S_{\beta} \in \vec{S}$ remains a Suslin tree in $W' [\forceP]$.
\end{lemma}

\begin{lemma}\label{nounwantedcodes}
If $\forceP \in W$ is allowable, $\forceP=(\forceP_{\beta} \, : \, \beta < \delta)$, $G \subset \forceP$ is generic over $W$ and $\{ x_{\beta} \, : \, \beta < \delta\}$ is the set of reals which are coded by $\forceP$. Let $\Phi(v_0)$ be the distinguished formula from above. Then
in $W[G]$, the set of reals which satisfy $\Phi(v_0)$ is exactly 
$\{ x_{\beta} \, : \, \beta < \delta\}$.
\end{lemma}

\begin{proof}
Let $G$ be $\forceP$ generic over $W$, we work in $W[G]$. Let $g= (b_{\beta} \, : \, {\beta} < \delta)$ be the set of the $\delta$-many $\omega_1$-branches through $S_{-1}$ we use to determine the coding areas in  the factors of $\forceP$.  Note that the family $\{b_{\beta} \,: \, \beta < \delta \}$ forms an almost disjoint family of subsets of $\omega_1$, that is two distinct members of  $\{b_{\beta} \,: \, \beta < \delta \}$ always have countable intersection. Thus there is an $\eta < \omega_1$ such that for arbitrary distinct $\beta_1$, $\beta_2 < \delta$,  $\eta> b_{\beta_1}\cap b_{\beta_2}$.  


We assume for a contradiction, that there is a real $(x,y,m)$ which satisfies $\Phi ((x,y,m))$ but $(x,y,m) \notin \{x_{\beta} \mid \beta < \delta \}$.
So there is a real $r$ such that $( {\ast} {\ast} {\ast})$ is true for $r$ and $(x,y,m)$. By the discussion above, $r$ carries enough information so that it codes a coding area $b$ and for every $\gamma \in b$ the following holds true.
\begin{align*}
n \in (x,y,m) \Rightarrow L[r] \models  ``S_{\omega \gamma + 2n+1} \text{ has an $\omega_1$-branch}".
\end{align*}
and
\begin{align*}
n \notin (x,y,m) \Rightarrow L[r] \models ``S_{\omega \gamma + 2n} \text{ has an $\omega_1$-branch}".
\end{align*}


As $(x,y,m)$ is distinct from every $x_{\beta}$, there must be $\aleph_1$-many $\alpha > \eta$ and an $n \in \omega$ such that, without loss of generality,
\[L[r] \models ``S_{\omega \alpha + 2n+1} \text{ has an $\omega_1$-branch}",\]
yet for every $r_{\beta}$ such that $({\ast} {\ast} {\ast} )$ holds true for $r_{\beta} $ and $x_{\beta}$,
\[L[r_{\beta}] \models  ``S_{\omega \alpha + 2n+1} \text{ does not have an $\omega_1$-branch}".\]

We fix an $\alpha$ which is as above. We claim  that there is no real in $W[G]$ such that $W[G] \models L[r] \models ``S_{\omega \alpha + 2n+1}$ has an $\omega_1$-branch$"$, which will be the desired contradiction.


We show this by pulling the generic branch for the forcing $S_{\omega \alpha + 2n+1}$ out of the generic for the  forcing $\forceQ^0 \ast \forceQ^1 \times \forceQ^2 \ast \forceP$ over $L$ which produces $W[G]$. 
Recall that the $\forceP$-generic $G$ produces $\delta$-many cofinal branches $\{b_{\beta} \mid \beta < \delta\}$ through $S_{-1}$. We let $\vec{S} \cup \{S_{-1}\}
 \ast \vec{b} \ast H$ denote a $(\forceQ^0 \ast \forceQ^1 \ast \forceQ^2)$-generic filter over $L$. Working in $W[G]$ we consider 
 \begin{align*}
A:= &\{ \omega \gamma +2n \mid \exists \beta < \delta (\gamma \in b_{\beta} \land n \notin x_{\beta}) \} \cup \\& \{\omega \gamma + 2n+1 \mid \exists \beta< \delta (\gamma \in b_{\beta} \land n \in x_{\beta} )\}.
\end{align*}
Recall that any almost disjoint coding forcing $\mathbb{A} (Y)$ can be defined in an absolute way over any universe containing $Y\subset \omega_1$.
As $\omega\alpha +2n+1 \notin A$  we note that in
\[ W[G]= L[\vec{S} \cup \{ S_{-1} \} ] [\{ b_i \mid i < \omega_1] [H] [G]\]
the forcing $\forceP^G$, whose factors are just branches through $S_{-1}$ and some almost disjoint coding forcings determined by the sets of the $x_{\beta}$'s and $b_{\beta}$'s, can already be defined in the inner model
\[ L[\vec{S}] [ \{ b_i \mid i < \omega_1, i \ne \omega\alpha +2n+1 \} [H] [G] \]
as can be seen by induction on the length of the iteration $\forceP$, using the fact that
$S_{\omega \alpha+2n+1}$ is still Suslin in $L[\vec{S}]  [\{ b_i \mid i < \omega_1, i \ne \omega\alpha +2n+1 \} [H] [G] $.

As a consequence, we can exploit the properties of product forcings to rearrange and obtain the equality
\begin{align*}
L[\vec{S} ] [ \{ b_i \mid i < \omega_1 \} ]& [H] [G]= \\& L[\vec{S}]  [ \{ b_i \mid i < \omega_1, i \ne \omega\alpha +2n+1 \} ] [H] [G] [b_{\omega\alpha +2n+1}] 
\end{align*}




To ease notation we let $B= \{ i <\omega_1 \mid i \ne \omega \alpha +2n +1 \}$
We claim now that $S_{\omega\alpha +2n+1}$ is a Suslin tree in $L[\vec{S} ] [\{b_{\beta} \mid \beta \in B \} ] [H] [G]$. This is clear by lemma \ref{PullingOutATree}.


Note now that, as forcing with $S_{\alpha+2n+1}$ is $\omega$-distributive, the reals will not change when passing from $L[\vec{S} ] [b_{\beta} \mid \beta \in B] [H] [G]$ to
$L[\vec{S} ] [b_{\beta} \mid \beta \in B] [H] [G] [b_{\alpha}]=W[G]$. 
 But this implies that 
\[L[\vec{S} ] [b_{\beta} \mid \beta \in B] [H] [G] \models \lnot \exists r L[r] \models `` S_{\alpha+2n+1} \text{ has an $\omega_1$-branch}" \]
as the existence of an $\omega_1$-branch through $S_{\alpha+2n+1}$ in the inner model $L[r]$ would imply the existence of such a branch in $L[\vec{S} ] [b_{\beta} \mid \beta \in B] [H] [G]$. Further, as no new reals appear when passing to $W[G]$ we also get 
\[W[G] \models \lnot \exists r L[r] \models `` S_{\alpha+2n+1} \text{ has an $\omega_1$-branch}". \]
This is the desired contradiction and as $G$ was an arbitrary generic filter the lemma is proven.

\end{proof}

\end{document}
